\documentclass[a4paper,11pt]{article}
\usepackage{url}
\usepackage{hyperref}
\usepackage{graphicx}
\usepackage{amsmath}
\usepackage{natbib}
\usepackage[margin=1in]{geometry}

\title{Community Detection: Topological vs. Embedding-Based Approaches Using the Markov Cluster Algorithm}
\author{Annafabia Gai, Vincenzo Sammartino, Riccardo Zanatta}
\date{18 December 2025}

\begin{document}

\maketitle

\section{What changed from the proposal}

\subsection{LINE embedding}
On the proposal we didn't discuss the method to reconstruct the graph from the embeddings. On this report (see following section) we defined in detail the methods for reconstruction and the intended experiments.

\subsection{Graph reconstruction}
On the proposal we stated that we intended on doing tests both on LINE and node2vec embeddings. It has been proven that the node2vec embeddings consistently outperform LINE embeddings on numerous tasks and in real applications are more commonly used. Given node2vec's superiority in comparison to LINE and the unavailability of LINE methods on updated Python libraries, we decided to drop LINE embeddings and work on the project only using node2vec.


\section{Addressing questions and requests}

\subsection{Graph reconstruction}
Starting from an input graph, we obtained node embeddings using the node2vec method. Next, employing the decoder component of the graph representation learning framework, we reconstructed the graph topology. To do so, we calculated pairwise similarities between embedding vectors and reconstructed the graph using k-Nearest Neighbor (k-NN). For every node $i$, the algorithm identifies the top $k$ most similar nodes (based on the similarity measure of choice) and creates edges from $i$ to them. Finally, we added on the reconstructed graph as edge weights the computed similarity between the nodes.
This methodology aligns with the framework proposed in a previous study (Yip et al., 2023), which assesses embedding quality explicitly through graph reconstruction. While RESTORE utilizes global ranking metrics (Precision@k), it relies on pairwise similarity scores derived from the embeddings specifically utilizing normalized dot products to measure node proximity. Similarly, large-scale industry recommender systems (e.g., Alibaba, PinSage) and similarity search infrastructure (e.g., Faiss) rely on retrieving nearest neighbors in vector space to approximate graph edges, validating k-NN as a robust standard for reconstruction.
Although node2vec implicitly optimizes a dot product, given that the method is based on Skip-Gram, the resulting vectors often exhibit magnitude biases correlated with node degree. As noted in the previous literature (PinSage and DeepWalking Backwards), applying normalization stabilizes the retrieval process. Therefore, we utilize cosine similarity for reconstruction, to prioritize structural alignment over magnitude, effectively normalizing the dot product to a more stable metric.
\newline
**cosa che mi è venuta in mente che secondo me è importante**
Using k-NN to reconstruct the graph has a strong effect on the graph connectivity and structure, introducing inductive bias by regularizing the degree distribution. Hubs are destroyed and the density landscape is flattened, leading to the breaking of large, messy communities into smaller, more uniform ones. Assigning a set value $k$ of edges to every node makes so that nodes with high connectivity in the real graph and nodes with low connectivity appear to be more similar in comparison to the original graph, even if the edge weights will be substantially different. In the context of community detection this may lead to less efficient community identification. Therefore, since the cited literature doesn't cover community detection in specific, to further evaluate graph reconstruction techniques in our context of study, we expanded our evaluation to include alternative reconstruction strategies:
\begin{itemize}
    \item Dot product + k-NN: to test if preserving magnitude information via raw dot product better retains hub prominence compared to cosine similarity.
    \item Cosine similarity + Thresholding: To test if a global similarity cutoff (allowing variable node degrees) preserves the natural density gradients better than a fixed $k$, despite the risk of disconnecting the graph.
    \item Dot Product + Thresholding: To evaluate the reconstruction preserving both magnitude and variable degree.
\end{itemize}

-> siccome Vandin dice (giustamente) che la ricostruzione del grafo è uno step cruciale del progetto, forse è il caso di provare a testare più metodi di ricostruzione. Questo potrebbe anche 'giustificare' il drop di LINE per favorire un'esplorazione più completa della situazione. Alla fine, una volta che è scritto il codice (avendo già fatto cosine + k-NN per scrivere le altre robe ci vogliono 2 secondi) basta runnare la roba un paio di volte in più e tirare giù le metriche.   

- the "Motivation" section of the report is somehow confusing. In fact, the title mentions only MCL, but then the Leiden algorithm is considered for topological clusterings.

- given your overall goal, there are three additional methods that you may want to include: MCL on the original graph; Leiden on the graph obtained from embeddings (you need to build such a graph for MCL, so you can use it for Leiden as well); k-means on the embeddings (with no MCL/Leiden)

- how large are the datasets that you will consider?

\section*{Dataset}
In the previous report, we lack in describing in detailed the dataset that we will imply. In particular, for this project we will use LFR (Lancichinetti–Fortunato–Radicchi) benchmark dataset to generate a
synthetic network used to evaluate the community detection algorithms used in our experiments. We decided to rely on a synthetic benchmark since it provides a groundtruth and a 
fair baseline that allows us to analyze how different parameters, conditions and implementation choices affect the quality of the detected community.

Indeed, the LFR benchmark offers a high degree of flexibility, as several hyperparameters can be tuned to generate networks with different structural properties.
For instance, we will focus on the μ parameter that controls how strongly communities are defined: low value μ defines well-separated communities and high level of μ weak or overlapping communities.

This parameter is crucial for our experiments because it reflects the performance limits of MCL and Leiden algorithms. 
As reported in the literature,MCL is known to perform poorly when communities are not well defined. For this reason, we investigate whether the use of node embeddings 
can improve the quality of the detected communities for both MCL and Leiden under different values of μ: from 0.2 and 0.3.

Moreover, according to other studies, a common choice for the network size in LFR-based evaluations is around 1000 nodes. Thus we will conduct our experiments on two kind of graphs: 500 nodes 
based graph and a 1000 nodes graph.
This allows us to check how the combination of the different value of μ and graph size affect the quality of community detection performance, and how
embeddings helps the algorithms to perform a better clustering.



\section*{Intended experiments}
As mentioned above, our aim is to explore how embeddings enhance MCL and Leiden under harsh networks condition for them to work nicely. 
Indeed we will manipulate the size N of the graph and the μ parameter to create different networking scenarios and check how the algorithms behave in the presence and not of embeddings.
We will perform the following steps, considering two different value for N and for μ:
First step explore a confortable scenario of network graph, where the number of nodes is not high and the communities are not strongly overlapping: N=500 and μ=0.2
- MCL without embeddings and with embeddings.
- Leiden without embeddings and with embeddings.

Then we will increase the value of N to 1000 and leaving μ=0.2; and computing:
- MCL without embeddings and with embeddings.
- Leiden without embeddings and with embeddings.

In this step we expect to see very different results from the previous one since MCL would start to provide worst results as N increases.

Eventually we perform the last experiment: N=1000 and increasing the overlap of the communities to μ=0.3; and computing:
- MCL without embeddings and with embeddings.
- Leiden without embeddings and with embeddings.

Below we provide a breif description for each algorithm:
-MCL




-Leiden:
The Leiden method is a community detection algorithm that optimizes modularity and through iteratively steps performs local node movements, 
refines the partition to remove disconnected components, and aggregates the graph into a smaller network.
However, like other modularity-based methods, Leiden suffers from the resolution limit, which may prevent the detection of very small communities, 
and its performance can degrade in networks with weak or highly overlapping community structures. 



\newpage

\section*{Contributions of Each Member}



\section*{AI Tools Usage Declaration}



\end{document}


